% TOP_PROBLEM: {"id": 3092, "title": "Chromatic number of associahedron", "category_id": 6, "category": {"id": 6, "name": "geometry", "display_name": "Geometry", "description": "Euclidean and non-Euclidean geometry, geometric structures.", "slug": "geometry", "order_index": 6, "created_at": "2026-07-31T15:26:25.671Z"}, "status": "open", "problem_number": "OPG-59984", "set_id": 12}
% FIRST_POSTED: 2026-10-01
% ATTEMPT_STATUS: unresolved
% COMPLETION: 5
% MODEL: GPT 6 Astra Ultra
% UnsolvedMath ID 3092; source catalogue code OPG-59984
% !TeX program = lualatex
% Research attempt, 2026-10-01. Canonical statement preserved verbatim.
\documentclass[11pt,a4paper]{article}
\usepackage[margin=25mm]{geometry}
\usepackage{fontspec}
\setmainfont{lmroman10-regular.otf}[BoldFont=lmroman10-bold.otf,
 ItalicFont=lmroman10-italic.otf,BoldItalicFont=lmroman10-bolditalic.otf]
\usepackage{amsmath,amssymb,mathtools}
\usepackage{unicode-math}
\setmathfont{latinmodern-math.otf}
\AtBeginDocument{\let\setminus\smallsetminus}
\usepackage{xurl,hyperref}
\hypersetup{colorlinks=true,urlcolor=blue}
\setcounter{secnumdepth}{0}
\setlength{\parindent}{0pt}
\setlength{\parskip}{5pt}
\setlength{\emergencystretch}{3em}
\begin{document}
\section*{Chromatic number of associahedron}\label{3092}
{\small UnsolvedMath ID 3092 · OPG-59984 · Geometry · OpenGarden}
\subsection{Definitions and mathematical statement}
For an integer $n\ge3$, fix a convex polygon with $n$ cyclically labelled
vertices. A triangulation is a maximal family of pairwise noncrossing
internal diagonals, or equivalently a subdivision into triangles using
only the polygon's vertices. It has $n-3$ internal diagonals.
Each internal diagonal borders two triangles whose union is a convex
quadrilateral. A flip removes that diagonal and inserts the other diagonal
of the quadrilateral.

Let $\mathcal A_n$ be the finite simple graph with triangulations as its
vertices and an edge between two triangulations exactly when one flip
transforms one into the other. This is the one-skeleton of the associahedron
associated with the $n$-gon; its polytope dimension is $n-3$.
A proper colouring of $\mathcal A_n$ assigns colours to triangulations
so that any two related by a flip get different colours. Its chromatic
number $\chi(\mathcal A_n)$ is the least number of colours required.

The conjecture asks whether
\[
 \sup_{n\ge3}\chi(\mathcal A_n)=\infty.
\]
Explicitly, for every positive integer $C$, must there be a polygon size
$n$ for which no colouring with $C$ colours is proper? The polygon's
actual lengths and angles are irrelevant; only cyclic vertex order
enters the graph. The problem concerns ordinary vertex colouring of the
flip graph, not colouring the diagonals within one triangulation, and
not colouring faces of the polytope. It does not specify a particular
growth rate beyond unboundedness.
\subsection{Short English statement}
Is there no fixed number of colours that properly colours triangulations of every polygon when a single flip must change colour?
\subsection{Research attempt}
\paragraph{Literature check and attempted route.}
Addario-Berry, Reed, Scott and Wood [A1] prove a logarithmic upper bound for
associahedron colouring. Their upper bound does not imply the requested
unbounded lower bound. The original problem [S2] discusses a lower bound of
four for a particular polygon; this attempt does not improve it. Instead,
it tests a natural amplification idea: insert many independent pentagon flip
cycles and try to multiply their colouring requirements. The calculation below
shows exactly why that amplification fails. Sources were checked 1 October 2026.

\paragraph{An induced product of pentagon cycles.}
For $r\ge1$ a convex $(3r+2)$-gon can be dissected into $r$ pentagons by
$r-1$ noncrossing diagonals: start with a pentagon and repeatedly glue another
pentagon along a boundary edge, thereby adding three vertices. Fix the
separating diagonals and consider the triangulations containing all of them.
A triangulation of each pentagonal cell can be chosen independently, and each
cell has five triangulations. Its flip graph is a $5$-cycle: there are two
flippable diagonals per triangulation and the five triangulations are connected
by successive flips.
A flip between two triangulations that both retain every separating diagonal
must occur inside one cell. Conversely every flip in one cell gives such an
edge. Hence the induced subgraph is exactly
\[
 X_r=\underbrace{C_5\mathbin\Box\cdots\mathbin\Box C_5}_{r\text{ factors}}.
\]
The word ``induced'' matters: flipping a separator leaves this vertex set
and cannot create an extra edge between two of its retained vertices.

\paragraph{An explicit three-colouring for every number of cells.}
Label the vertices of each $C_5$ cyclically by $0,1,2,3,4$ and use colours
$c=(0,1,0,1,2)$ in $\mathbb Z/3\mathbb Z$. Adjacent entries of this cyclic
list differ, including the final pair $2,0$. Define
\[
 C(t_1,\ldots,t_r)=\sum_{j=1}^r c(t_j)\pmod3.
\]
An edge of $X_r$ changes exactly one coordinate along a cycle edge, so its
colour difference is a nonzero element of $\mathbb Z/3\mathbb Z$. Thus $C$
is proper. Holding all but one coordinate fixed gives an odd cycle inside
$X_r$, proving that two colours are impossible. We have proved
\[
 |V(X_r)|=5^r,\qquad \deg(X_r)=2r,\qquad \chi(X_r)=3.
\]
Exponential growth in the number of triangulations and unbounded degree in
this family therefore produce no increasing chromatic obstruction.

\paragraph{The general limitation of independent-cell arguments.}
Suppose fixed separators divide a polygon into cells of sizes $n_j$, and
suppose each cell flip graph has a colouring into $k$ colours. Regard all
these colours as residues modulo the same $k$. Their sum, exactly as above,
properly colours the product graph. Each factor is also an induced subgraph,
obtained by fixing every other cell. It follows that the chromatic number of
this product is precisely the maximum of the factor chromatic numbers.
Thus independent repetition cannot amplify any fixed cell obstruction,
including one that requires four colours. Multiplying chromatic numbers would
be the wrong graph-product rule for these single-cell flips.

\paragraph{Verification and first remaining gap.}
A finite check enumerates all product states for $r=1,2,3,4$ and every
coordinate flip, verifying the formula above on $5,25,125,625$ vertices.
The proof covers all $r$. For $r=1$ it recovers the familiar pentagon, and
for $r=2$ it checks that commuting flips form squares, not extra diagonal
edges. To derive an unbounded lower bound, a construction must exploit flips
that interact across cells or a different obstruction to bounded colouring.
No such obstruction is established by this attempt; neither product size nor
degree supplies it.

\paragraph{Reproducibility.}
The finite calculations above are implemented in
\texttt{scripts/check\_attempt\_batch\_2026\_10\_01.py}; run it with Python 3
from the repository. It uses exact integer arithmetic and no external packages.

\paragraph{Final status.}
\textbf{UNRESOLVED.} The independent-pentagon amplification strategy is ruled
out by an explicit colouring, while the full associahedron conjecture is not
settled. These product-colouring facts are elementary and no novelty claim
is made. Completion estimate: 5\%.

\subsection{Sources}
\begin{itemize}
\item[S1] ULAM.AI and the UnsolvedMath Contributors, supplied problems.json, record 3092 (OPG-59984), OpenGarden collection. Catalogue identity and source transcription; CC BY 4.0.
\url{https://huggingface.co/datasets/ulamai/UnsolvedMath}.
\item[S2] Open Problem Garden, Chromatic number of associahedron. Original problem and discussion; the mathematical formulation below is expanded from the supplied source record.
\url{http://www.openproblemgarden.org/op/chromatic_number_of_associahedron}.
\item[A1] L. Addario-Berry, B. Reed, A. Scott and D. R. Wood, \emph{A logarithmic bound for the chromatic number of the associahedron}, arXiv:1811.08972. Primary abstract checked 1 October 2026. \url{https://arxiv.org/abs/1811.08972}.
\end{itemize}
\end{document}
